\frontchapter{Table des matières}
\begingroup\small\origpage{14}{总目录}
{\small\sffamily\color{BellaGray}以下页码为底本印刷页码，非本重排稿页码。}
\newcommand{\oldtoc}[3]{\par\noindent\hangindent=22mm\hangafter=1\makebox[22mm][l]{Chapitre #1}#2\nobreak\dotfill\nobreak\mbox{#3}\par}
\lecturetitle{CHAPITRES LIMINAIRES}
\oldtoc{1}{Existe-t-il une langue originelle ?}{002}
\oldtoc{2}{Le langage est-il le propre de l’homme ?}{024}
\lecturetitle{PARTIE 1 : LA GENÈSE D’UNE NOUVELLE SCIENCE}
\oldtoc{3}{De la grammaire à la linguistique : un art devenu science}{044}
\oldtoc{4}{Logique et linguistique : les deux sciences du logos}{071}
\oldtoc{5}{Linguistique comparative : la diversité des langues}{117}
\oldtoc{6}{Théories du signe : Peirce et Saussure}{150}
\lecturetitle{PARTIE 2 : LES SIX DOMAINES DES SCIENCES DU LANGAGE}
\oldtoc{7}{Phonétique articulatoire}{189}
\oldtoc{8}{Phonologie}{204}
\oldtoc{9}{Morphologie}{225}
\oldtoc{10}{Syntaxe}{270}
\oldtoc{11}{Sémantique}{304}
\oldtoc{12}{Pragmatique}{336}
\lecturetitle{PARTIE 3 : LES DÉVELOPPEMENT RÉCENTS DE LA LINGUISTIQUE\ednote{原文 « LES DÉVELOPPEMENT RÉCENTS » 中名词缺复数词尾，应为 « LES DÉVELOPPEMENTS RÉCENTS »。}}
\oldtoc{13}{Sociolinguistique : la prise en compte l’extralinguistique\ednote{原文如此；应为 « la prise en compte de l’extralinguistique »，在 « compte » 后补介词 « de »。}}{374}
\oldtoc{14}{Grammaire générative et révolution chomskyenne}{411}
\oldtoc{15}{Biolinguistique et psycholinguistique : le langage dans la tête}{456}
\oldtoc{16}{Langage et cognition : la linguistique cognitive}{486}
\vspace{1em}\noindent\textbf{BIBLIOGRAPHIE}\dotfill 526

\par\endgroup
